% ============================================================
%  ch10  第四讲 §4.1 单点单值化、§4.2 线丛与共形块
%  原文行 2199–2292
% ============================================================

\section{第四讲：对 $G$-丛模空间的应用}\label{Lecture IV}

\subsection{单点单值化}\label{one point uniformization}

设 $X$ 是 $k$ 上的光滑射影曲线，$\underline G$ 是 $X$ 上逐纤维连通的平滑仿射群概形。则 $\underline G$-主丛的模空间 $\Bun_{\underline G}$ 是 $k$ 上局部有限型的代数叠（例如见 \cite[Proposition 1]{He}）。

\begin{ex}
设 $\underline G|_{X-\{x_1,\ldots,x_n\}}\cong G\times (X-\{x_1,\ldots,x_n\})$，且在每个 $x_i$ 处 $\underline G|_{D_{x_i}}$ 是适当的抛孔群概形，则 $\Bun_{\underline G}$ 就是 $G$ 的抛物丛的模叠（抛物结构放在 $\{x_1,\ldots,x_n\}$ 处）。
\end{ex}

由 $\Gr_{\underline G,x}$ 的定义（参见 \eqref{aff glob}），有自然态射
\begin{equation}\label{one pt uni map}
u_x: \Gr_{\underline G,x}\to \Bun_{\underline G}.
\end{equation}
设 $X^*=X-\{x\}$，以 $\underline G^{X^*}$ 记截面空间，即
\begin{equation}\label{Gout}
\underline G^{X^*}(R)=\Gamma(X_R^*,\underline G).
\end{equation}

\begin{lem}\label{rep of Gout}
空间 $\underline G^{X^*}$ 可被一个群归纳概形表示。
\end{lem}

\begin{proof}
仍可化归到 $\underline G=\GL_n$，此时容易（例如见 \cite[(1.9)]{BL}）。
\end{proof}

通过限制到 $D_x^*\subset X^*$，可把 $\underline G^{X^*}$ 看作 $L\underline G_x$ 的子群。于是 $\underline G^{X^*}$ 沿 $u_x$ 的纤维传递作用，因此 $u_x$ 分解为
\begin{equation}\label{1pt uniform}
[\underline G^{X^*}\backslash\Gr_{\underline G,x}]\to \Bun_{\underline G}.
\end{equation}

该映射未必满射，$G=\bG_m$ 的情形即可见（开曲线 $X^*$ 上的线丛未必平凡）。但若 $\underline G=G\otimes X$ 半单，则它满射。事实上还有下述更强的结果。

\begin{thm}
设 $\underline G$ 是 $X$ 上逐纤维连通的平滑仿射群概形，其在有理函数域 $k(X)$ 上的通用纤维半单。则 \eqref{1pt uniform} 是同构。
\end{thm}

在 $k$-点层面，它说的是
\[\Bun_{\underline G}(k)= \Gamma\backslash G(F_x)/G(\mO_x),\]
其中 $\Gamma=\underline G^{X^*}(k)=\Gamma(X^*,\underline G)$。这归功于 Harder。$\SL_n$ 情形归功于 Beauville--Laszlo（见 \cite{BL}）。对常值半单群概形，这归功于 Drinfeld--Simpson（见 \cite{DS}）；\emph{同上}证明了：若 $p\nmid|\pi_1(G)|$，则 $\underline G^{X^*}$-主丛 $\Gr_{\underline G}\to \Bun_{\underline G}$ 在 \'etale 局部平凡。一般情形由 Heinloth 证明（见 \cite{He}）。因为有这个定理，可把 \eqref{one pt uni map} 看作单值化映射。

\begin{ex}
用最简单的例子 $G=\SL_2$ 说明该定理。此时 $\Bun_G$ 分类 $X$ 上行列式平凡的二阶向量丛。下面解释为什么这样的丛 $\mE$ 在 $X^*$ 上都变平凡。首先存在（足够负的）线丛 $\mL$ 使有短正合列 $0\to \mL\to \mE\to \mL^{-1}\to 0$。由于 $X^*$ 仿射，$\on{H}^1(X^*,\mL^{-2})=0$，故 $\mE|_{X^*}\cong (\mL\oplus\mL^{-1})|_{X^*}$。其次，可取足够大的整数 $n$，以及 $\mL(nx)$ 的整体截面 $s_1$ 与 $\mL^{-1}(nx)$ 的整体截面 $s_2$，使 $s_1$ 与 $s_2$ 在 $X^*$ 上无公共零点。令 $s=(s_1,s_2)|_{X^*}:\mO_{X^*}\to (\mL\oplus\mL^{-1})|_{X^*}$。则 $s$ 在 $X^*$ 上处处非零，于是有 $0\to \mO_{X^*}\to (\mL\oplus\mL^{-1})|_{X^*}\to \mO_{X^*}\to 0$。再由 $\on{H}^1(X^*,\mO_{X^*})=0$，得 $\mE|_{X^*}\cong (\mL\oplus\mL^{-1})|_{X^*}\cong \mO_{X^*}^2$。
\end{ex}

为简单起见，本小节余下部分假设 $\underline G=G\times X$ 常值，并把 $\underline G^{X^*}$ 记作 $G^{X^*}$。

\begin{cor}
有典范同构 $c:\pi_0(\Bun_G)\cong \pi_1(G)$。特别地，若 $G$ 单连通，则 $\Bun_G$ 连通。
\end{cor}

\begin{ex}
当 $G=\GL_n$ 时，映射 $\Bun_G\to \pi_0(\Bun_G)\cong\pi_1(G)\cong \bZ$ 由给 $X$ 上的向量丛配上其 Chern 类给出。
\end{ex}

\begin{proof}
要点是证明 $G^{X^*}$ 含于 $LG$ 的单位连通分支，于是 $\pi_0(\Bun_G)\cong \pi_0(\Gr_{x})$，再用定理 \ref{components}。该断言之证明见 \cite[(1.2)]{BLS}。
\end{proof}

\begin{rmk}
(i) 注意 $G^{X^*}$ 本身未必连通。事实上 \cite{BLS} 证明了 $\pi_0(G^{X^*})\simeq H^1(X^*,\pi_1(G))$（在 $\cha\ k=0$ 时）。

(ii) 该结果对非常值群概形的推广见 \cite{He}。
\end{rmk}

\subsection{线丛与共形块}

单值化定理对研究 $\Bun_G$ 上的线丛尤为有用。下面是主要定理。

\begin{thm}\label{uniformize line}
设 $G$ 单连通。则拉回诱导同构 $u_x^*:\Pic^e(\Bun_G)\to \Pic^e(\Gr_x)$。
\end{thm}

特别地，若 $G$ 单，则 $\Pic^e(\Bun_G)\cong\bZ$；$\Gr_G$ 的丰沛生成元 $\mO(1)$ 下降为 $\Pic^e(\Bun_G)$ 的生成元，仍记作 $\mO(1)$。

该映射的单射性容易说明：由单值化定理，$\Gr_G$ 上的刚性化线丛等同于 $\Gr_G$ 上 $G^{X^*}$-等变线丛，因此该推论等价于「任何特征 $G^{X^*}\to \bG_m$ 都平凡」；后者见 \cite[Corollary 5.2]{LS}。然而由于 $G^{X^*}$ 的几何仍很复杂，从单点单值化定理推出定理的满射部分并不容易。对典型群与 $G_2$，定理 \ref{uniformize line} 在 \cite{LS} 中通过显式构造拉回为 $\mO(1)$ 的 $\Bun_G$ 上线丛而得证（类似于注 \ref{generator} 的讨论）。完整定理先后由 Sorger \cite{So0}（$\cha\ k=0$，用到共形块空间的非消没结果）与 Faltings \cite[Theorem 17]{Fa}（一般情形）证明。下一小节将从另一个单值化定理推出它。

现在设 $k=\bC$（或任意特征零的代数闭域），$\mL$ 是 $\Bun_G$ 上的线丛。$\mL$ 的整体截面通常称为\textbf{非阿贝尔 theta 函数}（因为 $G=\bG_m$ 时就是经典 theta 函数）。设 $\frakg^{X^*}=\frakg\otimes\mO_X(X^*)$，这是 $G^{X^*}$ 的李代数。

\begin{lem}
Kac--Moody 中心扩张（即 \eqref{central extension} 的李代数）在 $\frakg^{X^*}$ 上唯一分裂，因此 $\frakg^{X^*}$ 作为子代数嵌入 Kac--Moody 代数 $\frakg_\Ga$（参见 \S\ref{Kac-Moody}）。
\end{lem}

\begin{proof}
由注 \ref{2-cocycle} 的显式公式与留数定理，中心扩张在 $\frakg^{X^*}$ 上分裂；又因 $[\frakg,\frakg]=\frakg$，这样的分裂唯一。
\end{proof}

在物理中，可积表示 $L_m(\frakg_\Ga)$ 是 WZW 共形场论的真空态空间，而 $\Hom_{\frakg^{X^*}}(L_m(\frakg_\Ga),k)$ 称为\textbf{共形块}空间。结合上述单值化定理与 Borel--Weil 型定理 \ref{Borel-Weil}，得

\begin{cor}
设 $G$ 单且单连通。$\Bun_G$ 上 $m$ 级非阿贝尔 theta 函数空间同构于共形块的对偶，即
\[\Gamma(\Bun_G,\mO(m))^*\simeq L_m(\frakg_\Ga)/\frakg^{X^*} L_m(\frakg_\Ga),\]
其中 $\mO(m)=\mO(1)^{\otimes m}$。
\end{cor}

共形块空间的维数由所谓的 Verlinde 公式计算，由此也得到 $\Bun_G$ 上非阿贝尔 theta 函数空间的维数。例如

\begin{cor}
设 $G=\SL_n$。则
\[\dim \Gamma(\Bun_G,\mO(m))= \Big(\frac{n}{n+m}\Big)^g \sum_{\substack{S\subset \{1,\ldots,m+n\}\\|S|=n}}\prod_{\substack{s\in S\\ t\not\in S}}\left|2\sin\pi\frac{s-t}{m+n}\right|^{g-1}.\]
\end{cor}

特别地，右端是正整数！见 \cite[Corollary 8.6]{BL}。

\begin{cor}
设 $G$ 是 $A,D,E$ 型，则 $\dim \Ga(\Bun_G,\mO(1))=|Z(G)|^g$，其中 $Z(G)$ 是 $G$ 的中心，$g$ 是 $X$ 的亏格。
\end{cor}

特别地，若 $G=E_8$，则 $\Bun_G$ 上定义 $\mO(1)$ 的除子唯一。在几何上找出这个除子将非常有趣（此问题由 Faltings 提出，见 \cite{Fa2}）。

单点单值化定理的进一步讨论与应用见 \cite{BL, LS, BLS, So0, So, He}，另见 \cite{KNR, KN}。
